\chapter{Descripción de casos de uso}
\label{anexo:casos_de_uso}

\begin{table}[H]
	\centering
	\begin{tabular}{| m{0.2\textwidth} | m{0.2\textwidth} | m{0.2\textwidth} | m{0.2\textwidth} |}
		\hline
		\rowcolor[HTML]{C0C0C0} Caso de Uso & \multicolumn{2}{|m{0.43\textwidth}|}{registro} &  CU1\\ 
		\hline
		Actores & \multicolumn{3}{l|}{Usuario, Administrador} \\ 
		\hline
		Tipo & \multicolumn{3}{l|}{Obligatorio} \\ 
		\hline
		Referencias & \multicolumn{3}{l|}{\hyperref[req:RF01]{RF\_01}} \\ 
		\hline
		Precondición & \multicolumn{3}{l|}{El usuario no debe estar registrado en el sistema.} \\ 
		\hline
        Postcondición & \multicolumn{3}{l|}{El usuario queda registrado en el sistema.} \\ 
		\hline
		Autor & \multicolumn{3}{l|}{Sergio Muñoz} \\ 
		\hline
		Fecha & 28/02/2025 & Versión & 1.0 \\
		\hline
	\end{tabular}
\end{table}

\begin{table}[H]
	\centering
	\begin{tabular}{| m{0.2\textwidth} | m{0.2\textwidth} | m{0.2\textwidth} | m{0.2\textwidth} |}
		\hline
		Propósito & \multicolumn{3}{m{0.67\textwidth}|}{El usuario debe ser capaz de crearse una cuenta en el sistema.}   \\ 
		\hline
		Resumen & \multicolumn{3}{m{0.67\textwidth}|}{El usuario introduce los datos necesarios para crear una cuenta.} \\ 
		\hline
	\end{tabular}
\end{table}

\begin{table}[H]
	\centering
	\begin{tabular}{| m{0.03\textwidth} | m{0.37\textwidth} | m{0.03\textwidth} | m{0.37\textwidth} |}
		\hline
		\multicolumn{4}{|m{0.9\textwidth}|}{Curso normal}     \\ 
		\hline
		1. & Usuario: Introduce su nombre, apellidos, correo electrónico y contraseña. 
        & 2. & El sistema verifica que los campos no estén vacíos y que la información introducida sea válida. \\ 
		\hline
		3. & Usuario: Confirma el registro. & 4. & El sistema verifica que el correo electrónico no esté registrado previamente.   \\ 
		\hline
        5. & Usuario: Recibe un mensaje indicando que debe verificar su correo electrónico. & 6. & El sistema redirige al usuario a la pantalla de inicio de sesión. \\ 
        \hline
		\multicolumn{4}{|m{0.9\textwidth}|}{Cursos alternos}     \\
        \hline 
        2a. & \multicolumn{3}{l|}{Los campos están vacíos o la información no es válida.} \\ 
        \hline
        2b. & \multicolumn{3}{l|}{El correo electrónico ya está registrado.} \\ 
        \hline
        2c. & \multicolumn{3}{l|}{El sistema muestra un mensaje de error.} \\ 
        \hline
	\end{tabular}
	\caption{Caso de uso - Registro}
\end{table}

\newpage

\begin{table}[H]
    \centering
    \begin{tabular}{| m{0.2\textwidth} | m{0.2\textwidth} | m{0.2\textwidth} | m{0.2\textwidth} |}
        \hline
        \rowcolor[HTML]{C0C0C0} Caso de Uso & \multicolumn{2}{|m{0.43\textwidth}|}{iniciarSesion} &  CU2\\ 
        \hline
        Actores & \multicolumn{3}{l|}{Usuario} \\ 
        \hline
        Tipo & \multicolumn{3}{l|}{Obligatorio} \\ 
        \hline
        Referencias & \multicolumn{3}{l|}{\hyperref[req:RF02]{RF\_02}} \\ 
        \hline
        Precondición & \multicolumn{3}{l|}{El usuario debe estar registrado en el sistema.} \\ 
        \hline
        Postcondición & \multicolumn{3}{l|}{El usuario accede al sistema.} \\ 
        \hline
        Autor & \multicolumn{3}{l|}{Sergio Muñoz} \\ 
        \hline
        Fecha & 28/02/2025 & Versión & 1.0 \\
        \hline
    \end{tabular}
\end{table}

\begin{table}[H]
    \centering
    \begin{tabular}{| m{0.2\textwidth} | m{0.2\textwidth} | m{0.2\textwidth} | m{0.2\textwidth} |}
        \hline
        Propósito & \multicolumn{3}{m{0.67\textwidth}|}{El usuario debe ser capaz de iniciar sesión en el sistema.}   \\ 
        \hline
        Resumen & \multicolumn{3}{m{0.67\textwidth}|}{El usuario introduce su correo electrónico y contraseña para autenticarse.} \\ 
        \hline
    \end{tabular}
\end{table}



\begin{table}[H]
    \centering
    \begin{tabular}{| m{0.03\textwidth} | m{0.37\textwidth} | m{0.03\textwidth} | m{0.37\textwidth} |}
        \hline
        \multicolumn{4}{|m{0.9\textwidth}|}{Curso normal}     \\ 
        \hline
        1. & Usuario: Introduce su correo electrónico y contraseña. 
        & 2. & El sistema verifica que los campos no estén vacíos y que la información introducida sea válida. \\ 
        \hline
        3. & Usuario: Confirma el inicio de sesión. & 4. & El sistema verifica que el correo electrónico y la contraseña sean correctos.   \\ 
        \hline
        5. & Usuario: Accede a la pantalla principal de la aplicación. & 6. & El sistema redirige al usuario a la pantalla principal de la aplicación. \\ 
        \hline
        \multicolumn{4}{|m{0.9\textwidth}|}{Cursos alternos}     \\
        \hline 
        2a. & \multicolumn{3}{l|}{Los campos están vacíos o la información no es válida.} \\ 
        \hline
        2b. & \multicolumn{3}{l|}{El correo electrónico o la contraseña son incorrectos.} \\ 
        \hline
        2c. & \multicolumn{3}{l|}{El sistema muestra un mensaje de error.} \\ 
        \hline
    \end{tabular}
    \caption{Caso de uso - Iniciar Sesión}
\end{table}

\newpage

\begin{table}[H]
    \centering
    \begin{tabular}{| m{0.2\textwidth} | m{0.2\textwidth} | m{0.2\textwidth} | m{0.2\textwidth} |}
        \hline
        \rowcolor[HTML]{C0C0C0} Caso de Uso & \multicolumn{2}{|m{0.43\textwidth}|}{cerrarSesion} &  CU3\\ 
        \hline
        Actores & \multicolumn{3}{l|}{Usuario} \\ 
        \hline
        Tipo & \multicolumn{3}{l|}{Obligatorio} \\ 
        \hline
        Referencias & \multicolumn{3}{l|}{\hyperref[req:RF04]{RF\_04}} \\ 
        \hline
        Precondición & \multicolumn{3}{l|}{El usuario debe estar autenticado en el sistema.} \\ 
        \hline
        Postcondición & \multicolumn{3}{l|}{El usuario cierra su sesión en el sistema.} \\ 
        \hline
        Autor & \multicolumn{3}{l|}{Sergio Muñoz} \\ 
        \hline
        Fecha & 28/02/2025 & Versión & 1.0 \\ 
        \hline
    \end{tabular}
\end{table}

\begin{table}[H]
    \centering
    \begin{tabular}{| m{0.2\textwidth} | m{0.2\textwidth} | m{0.2\textwidth} | m{0.2\textwidth} |}
        \hline
        Propósito & \multicolumn{3}{m{0.67\textwidth}|}{El usuario debe ser capaz de cerrar su sesión en el sistema.} \\ 
        \hline
        Resumen & \multicolumn{3}{m{0.67\textwidth}|}{El usuario cierra su sesión para finalizar su uso del sistema.} \\ 
        \hline
    \end{tabular}
\end{table}

\begin{table}[H]
    \centering
    \begin{tabular}{| m{0.03\textwidth} | m{0.37\textwidth} | m{0.03\textwidth} | m{0.37\textwidth} |}
        \hline
        \multicolumn{4}{|m{0.9\textwidth}|}{Curso normal} \\ 
        \hline
        1. & Usuario: Pulsa el botón de cerrar sesión. 
        & 2. & El sistema muestra un mensaje de confirmación. \\ 
        \hline
        3. & Usuario: Confirma que desea cerrar sesión. 
        & 4. & El sistema cierra la sesión del usuario y redirige a la pantalla de inicio de sesión. \\ 
        \hline
    \end{tabular}
    \caption{Caso de uso - Cerrar Sesión}
\end{table}

\newpage

\begin{table}[H]
    \centering
    \begin{tabular}{| m{0.2\textwidth} | m{0.2\textwidth} | m{0.2\textwidth} | m{0.2\textwidth} |}
        \hline
        \rowcolor[HTML]{C0C0C0} Caso de Uso & \multicolumn{2}{|m{0.43\textwidth}|}{verUsuario} &  CU4\\ 
        \hline
        Actores & \multicolumn{3}{l|}{Usuario, Administrador} \\ 
        \hline
        Tipo & \multicolumn{3}{l|}{Obligatorio} \\ 
        \hline
        Referencias & \multicolumn{3}{l|}{\hyperref[req:RF05]{RF\_05}} \\ 
        \hline
        Precondición & \multicolumn{3}{l|}{El usuario debe estar autenticado en el sistema.} \\ 
        \hline
        Postcondición & \multicolumn{3}{l|}{El usuario visualiza la información de su perfil.} \\ 
        \hline
        Autor & \multicolumn{3}{l|}{Sergio Muñoz} \\ 
        \hline
        Fecha & 28/02/2025 & Versión & 1.0 \\ 
        \hline
    \end{tabular}
\end{table}

\begin{table}[H]
    \centering
    \begin{tabular}{| m{0.2\textwidth} | m{0.2\textwidth} | m{0.2\textwidth} | m{0.2\textwidth} |}
        \hline
        Propósito & \multicolumn{3}{m{0.67\textwidth}|}{El usuario debe ser capaz de ver la información de su perfil.} \\ 
        \hline
        Resumen & \multicolumn{3}{m{0.67\textwidth}|}{El usuario accede a la sección de perfil para visualizar su información personal.} \\ 
        \hline
    \end{tabular}
\end{table}

\begin{table}[H]
    \centering
    \begin{tabular}{| m{0.03\textwidth} | m{0.37\textwidth} | m{0.03\textwidth} | m{0.37\textwidth} |}
        \hline
        \multicolumn{4}{|m{0.9\textwidth}|}{Curso normal} \\ 
        \hline
        1. & Usuario: Accede a la sección de perfil. 
        & 2. & El sistema muestra la información del perfil del usuario. \\ 
        \hline
    \end{tabular}
    \caption{Caso de uso - Ver Usuario}
\end{table}

\newpage

\begin{table}[H]
    \centering
    \begin{tabular}{| m{0.2\textwidth} | m{0.2\textwidth} | m{0.2\textwidth} | m{0.2\textwidth} |}
        \hline
        \rowcolor[HTML]{C0C0C0} Caso de Uso & \multicolumn{2}{|m{0.43\textwidth}|}{editarUsuario} &  CU5\\ 
        \hline
        Actores & \multicolumn{3}{l|}{Usuario, Administrador} \\ 
        \hline
        Tipo & \multicolumn{3}{l|}{Obligatorio} \\ 
        \hline
        Referencias & \multicolumn{3}{l|}{\hyperref[req:RF06]{RF\_06}} \\ 
        \hline
        Precondición & \multicolumn{3}{l|}{El usuario debe estar autenticado en el sistema.} \\ 
        \hline
        Postcondición & \multicolumn{3}{l|}{El usuario actualiza su información en el sistema.} \\ 
        \hline
        Autor & \multicolumn{3}{l|}{Sergio Muñoz} \\ 
        \hline
        Fecha & 28/02/2025 & Versión & 1.0 \\ 
        \hline
    \end{tabular}
\end{table}

\begin{table}[H]
    \centering
    \begin{tabular}{| m{0.2\textwidth} | m{0.2\textwidth} | m{0.2\textwidth} | m{0.2\textwidth} |}
        \hline
        Propósito & \multicolumn{3}{m{0.67\textwidth}|}{El usuario debe ser capaz de editar la información de su perfil.} \\ 
        \hline
        Resumen & \multicolumn{3}{m{0.67\textwidth}|}{El usuario accede a la sección de perfil y modifica su información personal.} \\ 
        \hline
    \end{tabular}
\end{table}

\begin{table}[H]
    \centering
    \begin{tabular}{| m{0.03\textwidth} | m{0.37\textwidth} | m{0.03\textwidth} | m{0.37\textwidth} |}
        \hline
        \multicolumn{4}{|m{0.9\textwidth}|}{Curso normal} \\ 
        \hline
        1. & Usuario: Accede a la sección de perfil. 
        & 2. & El sistema muestra la información del perfil del usuario. \\ 
        \hline
        3. & Usuario: Modifica la información deseada. 
        & 4. & El sistema verifica que la información introducida sea válida. \\ 
        \hline
        5. & Usuario: Guarda los cambios. 
        & 6. & El sistema actualiza la información del usuario y muestra un mensaje de confirmación. \\ 
        \hline
        \multicolumn{4}{|m{0.9\textwidth}|}{Cursos alternos} \\ 
        \hline 
        4a. & \multicolumn{3}{l|}{La información introducida no es válida.} \\ 
        \hline
        4b. & \multicolumn{3}{l|}{El sistema muestra un mensaje de error.} \\ 
        \hline
    \end{tabular}
    \caption{Caso de uso - Editar Usuario}
\end{table}

\newpage

\begin{table}[H]
    \centering
    \begin{tabular}{| m{0.2\textwidth} | m{0.2\textwidth} | m{0.2\textwidth} | m{0.2\textwidth} |}
        \hline
        \rowcolor[HTML]{C0C0C0} Caso de Uso & \multicolumn{2}{|m{0.43\textwidth}|}{eliminarUsuario} &  CU6\\ 
        \hline
        Actores & \multicolumn{3}{l|}{Administrador} \\ 
        \hline
        Tipo & \multicolumn{3}{l|}{Opcional} \\ 
        \hline
        Referencias & \multicolumn{3}{l|}{} \\ 
        \hline
        Precondición & \multicolumn{3}{l|}{El usuario debe estar autenticado en el sistema.} \\ 
        \hline
        Postcondición & \multicolumn{3}{l|}{El usuario es eliminado del sistema.} \\ 
        \hline
        Autor & \multicolumn{3}{l|}{Sergio Muñoz} \\ 
        \hline
        Fecha & 28/02/2025 & Versión & 1.0 \\ 
        \hline
    \end{tabular}
\end{table}

\begin{table}[H]
    \centering
    \begin{tabular}{| m{0.2\textwidth} | m{0.2\textwidth} | m{0.2\textwidth} | m{0.2\textwidth} |}
        \hline
        Propósito & \multicolumn{3}{m{0.67\textwidth}|}{El administrador debe ser capaz de eliminar un usuario del sistema.} \\ 
        \hline
        Resumen & \multicolumn{3}{m{0.67\textwidth}|}{El administrador selecciona un usuario y lo elimina del sistema.} \\ 
        \hline
    \end{tabular}
\end{table}

\begin{table}[H]
    \centering
    \begin{tabular}{| m{0.03\textwidth} | m{0.37\textwidth} | m{0.03\textwidth} | m{0.37\textwidth} |}
        \hline
        \multicolumn{4}{|m{0.9\textwidth}|}{Curso normal} \\ 
        \hline
        1. & Administrador: Accede a la lista de usuarios. 
        & 2. & El sistema muestra la lista de usuarios registrados. \\ 
        \hline
        3. & Administrador: Selecciona el usuario a eliminar. 
        & 4. & El sistema muestra un mensaje de confirmación. \\ 
        \hline
        5. & Administrador: Confirma la eliminación del usuario. 
        & 6. & El sistema elimina al usuario y muestra un mensaje de confirmación. \\ 
        \hline
    \end{tabular}
    \caption{Caso de uso - Eliminar Usuario}
\end{table}